\newpage
\begin{center}
  \textbf{\large ВВЕДЕНИЕ}
\end{center}
\addcontentsline{toc}{chapter}{ВВЕДЕНИЕ}

Кавитация – активно исследуемая область благодаря своему фундаментальному значению и множеству практических применений.
В медицине она используется для неинвазивных операций на глазах, точечной доставки лекарств, разрушения почечных камней и безыгольных инъекций.
В промышленности кавитация применяется в ультразвуковой очистке, обработке поверхностей и повышении эффективности морских винтов и турбин.
Понимание динамики кавитационных пузырей важно для оптимизации этих процессов и снижения рисков.
В нашем исследовании изучается поведение последовательности лазерно-индуцированных кавитационных пузырей вблизи торца волокна с помощью компьютерного моделирования методом конечных объемов (FVM).
Работа включает анализ ударной волны, возникающей при коллапсе пузырей, и распределения параметров, таких как давление и скорость.
Полученные данные позволяют оценить воздействие на окружающие материалы и способствуют улучшению безопасности и эффективности кавитационных технологий.
